% GENERATED FILE. Edit canonical lessons, then run npm run generate:books.
% canonical-source-hash: fnv1a64:f1c60a9297d60766
% canonical-lessons: LA-C290-velim, LA-C290-si-esset, LA-R290-would, LA-W290-epistula

\chapter{Would}
\label{ch:la-would}
\hlchaptermodality{290}

\begin{chapteropening}
\textbf{By the end of this chapter:} \emph{I can say politely what I would like, and what would happen if something that is not true were true, and write a short letter that uses it.}

Complete LA-W290-epistula: I can say politely what I would like, and what would happen if something that is not true were true, and write a short letter that uses it.
\end{chapteropening}

\section[velim, possim, mālim]{velim — I would like}

\label{lesson:LA-C290-velim}



\begin{quote}
Say \textquotedblleft{}I want to see Rome\textquotedblright{}. (\emph{Rōmam vidēre volō.})

\emph{Volō} states a fact about you: I want. Said to a stranger it can sound blunt. Latin softens it by changing the verb's mood.
\end{quote}

\begin{grammarlens}[title={Grammar Lens: velim, the softer want}]
\begin{quote}
\textbf{Velim tē vidēre.} — I would like to see you.
\end{quote}

\textbf{velim} is the \textbf{subjunctive} of \emph{volō}: not \textquotedblleft{}I want\textquotedblright{}, but \textquotedblleft{}I would want, I would like\textquotedblright{}. All six persons:

\begin{quote}
\textbf{velim}, \textbf{velīs}, \textbf{velit}, \textbf{velīmus}, \textbf{velītis}, \textbf{velint}.
\end{quote}

The other wanting and able verbs do the same, with the same endings: \textbf{possim}, \textquotedblleft{}I could, I might be able\textquotedblright{}; \textbf{mālim}, \textquotedblleft{}I would rather\textquotedblright{}.

\begin{quote}
\textbf{Hoc facere possim.} — I could do this. \textbf{Mālim manēre.} — I would rather stay.
\end{quote}

Cicero writes \emph{velim scire}, \textquotedblleft{}I should like to know\textquotedblright{}, in one of his dialogues, and \emph{velim} again in a letter telling his friend Atticus how much he would like to see him (Lewis \& Short quote both).
\end{grammarlens}

\begin{culture}
\emph{Velim} is not past or future. It is a different way of holding the same present: the \textbf{indicative}, \emph{volō}, says what is; the \textbf{subjunctive}, \emph{velim}, says what would be, or could be. \emph{Sum} has one too, \emph{sim}, \textquotedblleft{}I may be\textquotedblright{}, and the next lesson puts \emph{sum}'s other subjunctive to work.
\end{culture}

\begin{cousinweb}
The \textbf{-i-} of \emph{velim}, \emph{possim} and \emph{sim} is usually explained as the mark of an old wishing mood, which Greek kept as a mood of its own and Latin folded into the subjunctive. Lewis \& Short do record an older \emph{siem}, \emph{siet} for \emph{sim}, \emph{sit} in Plautus; the link to the Greek mood is the account histories of Latin give, offered as one.
\end{cousinweb}

\subsection*{Your turn}
Say these aloud:

\begin{itemize}
\raggedright
  \item \emph{velim, velīs, velit} — then \emph{velīmus, velītis, velint}
  \item \emph{Velim tē vidēre.}
  \item \emph{Hoc facere possim.}
  \item \emph{Mālim manēre.}
  \item the blunt want, then the soft one — \emph{volō}, then \emph{velim}
  \item \emph{Sī vīs, tē vidēre velim.}
\end{itemize}

\begin{tcolorbox}[breakable,colback=teal!4,colframe=teal!35!black,title={Before you move on}]
Say \textquotedblleft{}I would like to see you\textquotedblright{}. (\emph{\textbf{Velim tē vidēre.}}) Say \textquotedblleft{}I could\textquotedblright{}. (\emph{\textbf{Possim.}}) Say \textquotedblleft{}I would rather stay\textquotedblright{}. (\emph{\textbf{Mālim manēre.}}) Is \emph{velim} a tense? (\textbf{No: a mood}, the subjunctive.)
\end{tcolorbox}
\section[sī hīc essēs, laetus essem]{sī hīc essēs — if it were so}

\label{lesson:LA-C290-si-esset}



\begin{quote}
Say \textquotedblleft{}if it rains, we stay in the farmhouse\textquotedblright{}. (\emph{Sī pluit, in vīllā manēmus.}) Then \textquotedblleft{}I would like to see you\textquotedblright{}. (\emph{Velim tē vidēre.})

The first may well come true. Now a condition you know is false, about right now.
\end{quote}

\begin{grammarlens}[title={Grammar Lens: the infinitive plus an ending}]
\begin{quote}
\textbf{Sī hīc essēs, laetus essem.} — If you were here, I would be happy.
\end{quote}

You are not here, and both halves say so by their mood. \textbf{essem} and \textbf{essēs} are the \textbf{imperfect subjunctive} of \emph{sum}, and they are made the easiest way Latin has: take the infinitive, \emph{esse}, and add the endings: \textbf{essem}, \textbf{essēs}, \textbf{esset}, \textbf{essēmus}, \textbf{essētis}, \textbf{essent}.

The other verbs follow the same recipe: \emph{habēre}, \textbf{habērem}, I would have; \emph{posse}, \textbf{possem}, I could; \emph{velle}, \textbf{vellem}, I would want.

\begin{quote}
\textbf{Sī pecūniam habērem, nāvem emerem.} — If I had money, I would buy a ship.
\end{quote}

\emph{Emerem} is \emph{emere} with \emph{-m}. In English the \textquotedblleft{}if\textquotedblright{} half sounds past, \textquotedblleft{}if I had\textquotedblright{}; Latin agrees, and calls this the imperfect.
\end{grammarlens}

\begin{culture}
Set the two side by side:

\begin{quote}
\textbf{Sī pluit, in vīllā manēmus.} — If it rains, we stay in. (It may.) \textbf{Sī nōn plueret, nāvigārēmus.} — If it were not raining, we would be sailing. (It is.)
\end{quote}

Plain verbs leave the question open. The subjunctive tells the listener the \textquotedblleft{}if\textquotedblright{} is false right now, and the speaker knows it.
\end{culture}

\subsection*{Your turn}
Say these aloud:

\begin{itemize}
\raggedright
  \item \emph{esse}, then \emph{essem, essēs, esset}
  \item \emph{Sī hīc essēs, laetus essem.}
  \item \emph{Sī pecūniam habērem, nāvem emerem.}
  \item \emph{Sī nōn plueret, nāvigārēmus.}
  \item one thing that is not true now, and what would follow if it were
\end{itemize}

\begin{tcolorbox}[breakable,colback=teal!4,colframe=teal!35!black,title={Before you move on}]
Say \textquotedblleft{}if you were here, I would be happy\textquotedblright{}. (\emph{\textbf{Sī hīc essēs, laetus essem.}}) What is \emph{essem} built on? (The \textbf{infinitive}, \emph{esse}.) Say \textquotedblleft{}I would have\textquotedblright{}. (\emph{\textbf{Habērem.}}) Which kind of condition says the \textquotedblleft{}if\textquotedblright{} is false? (The one with the \textbf{subjunctive}.)
\end{tcolorbox}
\section[(review)]{Review — what is, what may be, what would be}

\label{lesson:LA-R290-would}



\begin{quote}
Say \textquotedblleft{}I would like to see you\textquotedblright{}. (\emph{Velim tē vidēre.}) Then \textquotedblleft{}if you were here, I would be happy\textquotedblright{}. (\emph{Sī hīc essēs, laetus essem.})
\end{quote}

\begin{grammarlens}[title={What you've built}]
Two chapters, and three ways to hold a thought. \textbf{What is}, and what may well be: plain verbs, \emph{sī pluit, manēmus}. \textbf{What depends}: \emph{ex tempestāte pendet}. \textbf{What would be}: the subjunctive, \emph{velim} and \emph{possim} for the polite and the possible, \emph{essem} and \emph{habērem} for an \textquotedblleft{}if\textquotedblright{} that is not true now.

\begin{quote}
\emph{Sī pluit, in vīllā manēmus. Sed sī nōn plueret, nāvigārēmus, et laetus essem.}
\end{quote}

The first sentence leaves the weather open. The second knows it is raining, and wishes it were not.
\end{grammarlens}

\subsection*{Your turn}
Say these aloud:

\begin{itemize}
\raggedright
  \item \emph{Sī pecūniam habērem, nāvem emerem.}
  \item \emph{Hoc facere possim, sed mālim manēre.}
  \item \emph{In tuā potestāte est.} — then \emph{Spērō tē valēre.}
  \item \emph{Sī nōn plueret, nāvigārēmus.}
  \item one real condition and one unreal one about today
\end{itemize}

\begin{tcolorbox}[breakable,colback=teal!4,colframe=teal!35!black,title={Before you move on}]
Say \textquotedblleft{}I could\textquotedblright{}. (\emph{\textbf{Possim.}}) Say \textquotedblleft{}if I had money\textquotedblright{}. (\emph{\textbf{Sī pecūniam habērem.}}) In \emph{sī pluit, manēmus}, may it rain? (\textbf{Yes}, it is open.) In \emph{sī nōn plueret}, is it raining? (\textbf{Yes}: that is why the verb is subjunctive.) What is the \emph{-i-} of \emph{velim} usually said to be? (An \textbf{old wishing mood}.)
\end{tcolorbox}
\section[(epistula)]{A letter about a journey}

\label{lesson:LA-W290-epistula}



\begin{quote}
Your four lines with joiners held together. These six lines are an \textbf{epistula}, a letter, and they tell a story.

Pick one friend who will read it.
\end{quote}

\subsection*{What to know first}
A Roman letter opened with two names and two letters: \emph{Mārcus Iūliae S. D.}, Marcus to Julia, \textbf{S. D.}, \emph{salūtem dīcit}, \textquotedblleft{}sends greeting\textquotedblright{}. Your name goes first, in the nominative; your friend's in the dative.

Then, often, a real condition: \textbf{Sī valēs, bene est; ego valeō.} \textquotedblleft{}If you are well, good; I am well.\textquotedblright{} Seneca calls it an old custom for the first words of a letter (Lewis \& Short quote him), and it was common enough to be written as initials.

The middle tells the story: the scene in the imperfect, the events in the perfect, in order. The end looks forward, and the last word is \textbf{valē}.

\subsection*{Writing — connected composition}
On paper, write \textbf{six lines} to your friend about a journey, real or invented, one line for each of these:

\begin{enumerate}
\raggedright
  \item the opening: your name, your friend's, \emph{S. D.}, and \emph{sī valēs, bene est};
  \item the scene: when it was, with \emph{abhinc}, and the weather, with \emph{erat} or \emph{-bat};
  \item the first events, with \emph{prīmum} and \emph{deinde}: \emph{iter faciō}, a ship, the fare;
  \item what went wrong, with \emph{subitō} or \emph{quid accidit}, in the perfect;
  \item how you felt and why, with \emph{eram} and \emph{quia} or \emph{quod}, and one comparison with \emph{quam};
  \item one line forward, with \emph{sī … essēs, … essem}, or \emph{velim}, or \emph{spērō}, then \emph{valē}.
\end{enumerate}

Use only words you already write. The exercise is choosing the right past for each line.

\begin{tcolorbox}[breakable,colback=teal!4,colframe=teal!35!black,title={Before you move on}]
Now go through your lines as you went through the rainy day in your first story. Mark every verb \textbf{S} for scene or \textbf{E} for event. Line two should be all \textbf{S}; lines three and four mostly \textbf{E}; line five mostly \textbf{S}. A line about the weather marked \textbf{E} is the commonest slip: change \emph{pluit} to \emph{pluēbat}.

Then two endings to check. The perfect for \textquotedblleft{}they\textquotedblright{} ends in \textbf{-ērunt}; the imperfect has \textbf{-ba-} before its ending. And a feeling word agrees with whoever felt it: \emph{laeta eram} if a woman is writing.

Last, read line six on its own. If its \emph{sī} is about something that may be true, are both verbs plain? If it is about something false, are both subjunctive, \emph{essēs} and \emph{essem}? Then the letter looks forward, and it is finished.
\end{tcolorbox}
